% ============================================================================
% example_chapter.tex — CHAPTER RHYTHM gallery (content-style DNA).
% Not a box catalog:  a density ruler + 3 full short chapters that show the
% target cadence: one idea → one short box → numbered equation → short intuition.
%
%   bash ../scripts/build.sh example_chapter.tex
%   pdftoppm -png -r 140 example_chapter.pdf ch
%
% THIS is what to imitate when writing a study doc.  example_boxes.tex is only
% for picking a box *type* — not for chapter density.
% ============================================================================
% ============================================================================
% preamble.tex — tested LuaLaTeX preamble for Hebrew + math + figures PDFs (v3.23)
% Battle-tested: compiled an 80-page Hebrew quantum-computing summary AND a
% 197-page Hebrew computer-architecture anthology (TikZ diagrams, code listings,
% colored boxes) — both with 0 missing characters.
%
% Usage in your document:
%   % ============================================================================
% preamble.tex — tested LuaLaTeX preamble for Hebrew + math + figures PDFs (v3.23)
% Battle-tested: compiled an 80-page Hebrew quantum-computing summary AND a
% 197-page Hebrew computer-architecture anthology (TikZ diagrams, code listings,
% colored boxes) — both with 0 missing characters.
%
% Usage in your document:
%   % ============================================================================
% preamble.tex — tested LuaLaTeX preamble for Hebrew + math + figures PDFs (v3.23)
% Battle-tested: compiled an 80-page Hebrew quantum-computing summary AND a
% 197-page Hebrew computer-architecture anthology (TikZ diagrams, code listings,
% colored boxes) — both with 0 missing characters.
%
% Usage in your document:
%   \input{preamble.tex}     % (or paste these lines at the top)
%   \begin{document}  ...  \end{document}
%
% Compile with:  lualatex -interaction=nonstopmode doc.tex   (twice, for TOC)
% Requires the Hebrew fonts installed (run scripts/setup_fonts.sh first).
%
% SHORT doc (exam + solution): change `report` -> `article`.
% STUDY doc meant to be read/annotated: 12pt + \linespread{1.3} below are
%   deliberate readability choices (NOT padding) — keep or trim to taste.
% ============================================================================
\documentclass[12pt]{report}          % 12pt: readable for a study guide; 11pt is fine too
\usepackage[no-math]{fontspec}        % (1) ★ no-math — THE critical line (see recipe.md)
\usepackage[bidi=basic]{babel}        % (2) RTL handling built into the engine
\babelprovide[import,main]{hebrew}    % (3) Hebrew = main language, default RTL
\babelprovide[import]{english}        % (4) English secondary, for correct LTR runs
% Latin fallback registered BEFORE \babelfont so a lone English word in Hebrew
% prose does not tofu. Fonts resolved by FAMILY NAME (portable; no hard path).
\directlua{luaotfload.add_fallback("hebfb",{"DejaVuSerif:mode=node;"})}
% v3.22: Frank Ruhl Libre ships NO italic face. Instead of silently-upright
% \emph (v<=3.20) or bold-as-emphasis (v3.21), synthesize a slanted face with
% FakeSlant — real, VISIBLE italic emphasis with correct \emph nesting
% semantics (inner \emph toggles back upright). Verified empirically at
% 300dpi: slant clearly visible, 0 errors, 0 missing glyphs.
\babelfont{rm}[RawFeature={fallback=hebfb},
  ItalicFont={Frank Ruhl Libre}, ItalicFeatures={FakeSlant=0.22},
  BoldItalicFont={Frank Ruhl Libre Bold}, BoldItalicFeatures={FakeSlant=0.22}
  ]{Frank Ruhl Libre}
\babelfont{sf}[RawFeature={fallback=hebfb}]{Heebo}
\babelfont{tt}[RawFeature={fallback=hebfb}]{DejaVu Sans Mono}

\usepackage{amsmath,amssymb,mathtools}
\usepackage[margin=2.3cm,headheight=14pt]{geometry}
\linespread{1.3}                      % readable line spacing for study/annotation docs
\usepackage{xcolor}
\usepackage{mdframed}
\usepackage{listings}                 % code listings (assembly / C / pseudo-code)
\usepackage{tikz}
\usetikzlibrary{arrows.meta,positioning,calc,shapes.geometric,fit,backgrounds}
\usepackage{pgfplots}\pgfplotsset{compat=1.18}
\usepackage{quantikz}                 % quantum circuits (remove if unused)
\usepackage{booktabs}
\usepackage{longtable}
\usepackage{array}
\usepackage{float}                    % [H] = place table exactly here (reference doc: no drifting)
\usepackage{enumitem}
\usepackage{needspace}

% ============================================================================
% ★★ v3.4 — THE DIRECTION-ISOLATION TRIO (the heart of mixing Hebrew + Latin)
% Mental model:
%   \h{...}    Hebrew (RTL) island      — Hebrew text inside an LTR context
%                                          (TikZ nodes, the english `tikzpic`).
%   \en{...}   English (LTR) isolate    — Latin text / parentheses inside a
%                                          Hebrew (RTL) context. KEEPS PARENS
%                                          AND ORDER CORRECT.
%   \code{...} \en + monospace          — code inside Hebrew / inside a frametitle.
% CRITICAL: never wrap Latin in \h{} (it is for HEBREW). Inside a TikZ node,
% Latin in \h{} renders REVERSED (Datapath->htapataD, CPU->UPC) and parens
% mirror. Isolate every Latin/paren chunk with \en{} or \code{}. See
% references/figures-boxes-listings.md.
% ---- TEXT INSIDE MATH (v3.8) ----------------------------------------------
% Put ANY text that sits inside math in a \text{}. The ONLY thing that is
% invalid in math is a bare \h/\en (= \foreignlanguage) used DIRECTLY, outside
% \text -> "\rmfamily invalid in math mode". Inside \text:
%   Hebrew label           -> \text{...}            (bare Hebrew renders fine;
%                              \h is OPTIONAL/harmless: \text{\h{...}} also works
%                              and both render identically + scale in scripts)
%   plain Latin word        -> \text{...}            (no marker needed)
%   Latin WITH ( ) or [ ]   -> \text{\en{...}}       (\en forces LTR; a plain
%                              \text mirrors the parens in RTL -> "Spec (TNR(" )
%   f(Hebrew/text arg)      -> \text{f}(\text{...})  -- keep the parens at MATH
%                              level, NOT \text{f(arg)} (parens inside the box mirror)
%   code in math            -> \mathtt{...}  (or \texttt{}, math-guarded below)
% NEVER write \en{...} or \h{...} straight into $...$ / \[..\] (outside \text);
% always via \text{}. (\mbox{...} also works but does not scale in scripts.)
% ============================================================================
\newcommand{\h}[1]{\foreignlanguage{hebrew}{#1}}
\newcommand{\en}[1]{\foreignlanguage{english}{#1}}

% ============================================================================
% ★★ — fixes for defects that are INVISIBLE in the log and in the checks
% and only surface on a VISUAL RENDER of the PDF. (Each verified empirically.)
% ----------------------------------------------------------------------------
% v3.22: \emph is now REAL slanted emphasis (FakeSlant italic face defined on
% the \babelfont{rm} line above), replacing the v3.21 \emph->\textbf hack: it
% preserves \emph's toggle semantics and keeps bold free for its own role.
% Prefer bold-as-emphasis for your doc? one line: \renewcommand{\emph}[1]{\textbf{#1}}
% page number: force the digits LTR so they never reverse (18 -> 81) in the RTL
%   footer. That reversal is a NON-DETERMINISTIC bidi glitch (observed once in
%   ~80 pp); this LuaLaTeX build is not byte-reproducible, so the defensive
%   LTR wrap is the only reliable guard. Latin digits are always LTR anyway.
\renewcommand{\thepage}{\foreignlanguage{english}{\arabic{page}}}
% emergencystretch: Hebrew does not hyphenate, so an unbreakable LTR island
%   (\en{...} / an inline math run) landing at the margin can leave a small
%   (~2-13pt) overfull with no legal breakpoint. This rescues ONLY those
%   inline-prose tie-breaks; it is a ceiling used only on offending lines and
%   does NOT touch display-math / listing / table overflow (those still surface
%   in the log for the content-level overflow pass). Verified: baseline overfull
%   -> 0 with this set, with no change to pagination or underfull.
\setlength{\emergencystretch}{3em}
% em-dash: --- does NOT become — under Frank Ruhl (the font lacks tlig; verified
%   the width is identical with/without Ligatures=TeX, while the mechanism works
%   on Latin Modern). TYPE THE CHARACTER — (U+2014) DIRECTLY; likewise – (U+2013).
%   check_bidi_figures.py flags a --/--- that is adjacent to Hebrew.
% ============================================================================

% ---- STRUCTURAL DIRECTION-SAFETY for \texttt (v3.7) -------------------------
% A parenthesis, bracket, or $ INSIDE \texttt{} mirrors/migrates in RTL prose
% (MIPS 0($s0)->"(s0$)0", $rs->"rs$", arr[i]->"[arr[i") because \texttt sets the
% font but NOT the direction. Redefining \texttt to LTR-isolate its content fixes
% the WHOLE class at once (parens, brackets, $) so plain \texttt{lw $t0,0($s0)}
% is safe — no need to remember \code{}. Guarded by \ifmmode because
% \foreignlanguage is invalid in math (there, plain monospace is already LTR).
% NB: \texttt{Hebrew} stays tofu — the monospace font has no Hebrew anyway; put
% Hebrew in \h{}/normal text, never in \texttt{}. To disable, delete this block.
\let\origtexttt\texttt
\renewcommand{\texttt}[1]{\ifmmode\origtexttt{#1}\else\foreignlanguage{english}{\origtexttt{#1}}\fi}
\newcommand{\code}[1]{\ifmmode\origtexttt{#1}\else\foreignlanguage{english}{\origtexttt{#1}}\fi}

% ---- a thin decorative rule helper ----
\newcommand{\sep}{\par\vspace{2pt}\noindent\textcolor{cRule}{\rule{\linewidth}{0.6pt}}\par\vspace{4pt}}

% ---- colors ----
\definecolor{cBlue}{HTML}{1F4E79}   \definecolor{cBlueBg}{HTML}{EAF1F8}
\definecolor{cGreen}{HTML}{2E6B4F}  \definecolor{cGreenBg}{HTML}{EAF4EE}
\definecolor{cAmber}{HTML}{9C6A00}  \definecolor{cAmberBg}{HTML}{FBF3E2}
\definecolor{cPurple}{HTML}{5B2C6F} \definecolor{cPurpleBg}{HTML}{F3ECF7}
\definecolor{cTeal}{HTML}{0E6E6E}   \definecolor{cTealBg}{HTML}{E6F4F4}
\definecolor{cRed}{HTML}{8E2B2B}    \definecolor{cRedBg}{HTML}{FBEAEA}
\definecolor{cGray}{HTML}{555555}   \definecolor{cRule}{HTML}{1F4E79}

% ---- math operators ----
\DeclareMathOperator{\Tr}{Tr}
\DeclareMathOperator{\rank}{rank}
\DeclareMathOperator{\diag}{diag}
\DeclareMathOperator{\sgn}{sgn}

% ---- bra-ket (quantikz defines \ket/\bra/\braket; add the rest safely) ----
\providecommand{\ket}[1]{\left|#1\right\rangle}
\providecommand{\bra}[1]{\left\langle#1\right|}
\providecommand{\braket}[2]{\left\langle#1\middle|#2\right\rangle}
\newcommand{\dyad}[1]{\left|#1\right\rangle\!\left\langle#1\right|}
% v3.21: general outer product |a><b| and expectation value <A> — agents reach
% for these by reflex (esp. \expval); quantikz defines neither, so their absence
% threw "Undefined control sequence" in a real build. \providecommand = safe.
\providecommand{\ketbra}[2]{\left|#1\right\rangle\!\left\langle#2\right|}
\providecommand{\expval}[1]{\left\langle#1\right\rangle}
\newcommand{\abs}[1]{\left|#1\right|}
\newcommand{\norm}[1]{\left\lVert#1\right\rVert}
\newcommand{\half}{\tfrac12}
\newcommand{\R}{\mathbb{R}}\newcommand{\C}{\mathbb{C}}\newcommand{\Z}{\mathbb{Z}}
% identity operator. NB: \mathbb{1} silently renders a WRONG glyph (msbm has no
% blackboard digits). Use \mathbb{I} = 𝕀 (works with amssymb, no extra package).
% Do NOT "fix" this back to \mathbb{1}.
\newcommand{\id}{\mathbb{I}}

% ============================================================================
% Colored boxes — mdframed (NOT tcolorbox, which breaks under bidi).
% ★★ v3.4: innerleftmargin/innerrightmargin give code listings room so their
% frame/line-numbers do not spill out the box border (see listing styles below).
% ============================================================================
\newmdenv[linecolor=cBlue,backgroundcolor=cBlueBg,linewidth=1.1pt,
  frametitlebackgroundcolor=cBlue,frametitlefontcolor=white,
  innertopmargin=6pt,innerbottommargin=7pt,skipabove=9pt,skipbelow=9pt,
  innerleftmargin=9pt,innerrightmargin=9pt,roundcorner=3pt]{defbox}   % הגדרה
\newmdenv[linecolor=cGreen,backgroundcolor=cGreenBg,linewidth=1.1pt,
  frametitlebackgroundcolor=cGreen,frametitlefontcolor=white,
  innertopmargin=6pt,innerbottommargin=7pt,skipabove=9pt,skipbelow=9pt,
  innerleftmargin=9pt,innerrightmargin=9pt,roundcorner=3pt]{thmbox}   % משפט/טענה
\newmdenv[linecolor=cAmber,backgroundcolor=cAmberBg,linewidth=1.1pt,
  frametitlebackgroundcolor=cAmber,frametitlefontcolor=white,
  innertopmargin=6pt,innerbottommargin=7pt,skipabove=9pt,skipbelow=9pt,
  innerleftmargin=9pt,innerrightmargin=9pt,roundcorner=3pt]{notebox}  % הערה
\newmdenv[linecolor=cPurple,backgroundcolor=cPurpleBg,linewidth=1.1pt,
  frametitlebackgroundcolor=cPurple,frametitlefontcolor=white,
  innertopmargin=6pt,innerbottommargin=7pt,skipabove=9pt,skipbelow=9pt,
  innerleftmargin=9pt,innerrightmargin=9pt,roundcorner=3pt]{exbox}    % דוגמה
\newmdenv[linecolor=cRed,backgroundcolor=cRedBg,linewidth=1.1pt,
  frametitlebackgroundcolor=cRed,frametitlefontcolor=white,
  innertopmargin=6pt,innerbottommargin=7pt,skipabove=9pt,skipbelow=9pt,
  innerleftmargin=9pt,innerrightmargin=9pt,roundcorner=3pt]{warnbox}  % טעות נפוצה
\newmdenv[linecolor=cTeal,backgroundcolor=cTealBg,linewidth=1.1pt,
  frametitlebackgroundcolor=cTeal,frametitlefontcolor=white,
  innertopmargin=6pt,innerbottommargin=7pt,skipabove=9pt,skipbelow=9pt,
  innerleftmargin=9pt,innerrightmargin=9pt,roundcorner=3pt]{keybox}   % רעיון מרכזי

% ============================================================================
% Unified TABLE presentation: subtle uniform slate frame + auto-numbered
% caption rendered INSIDE the frame. Use:
%   \begin{tablebox}\centering\small  <tabular>  \tabcap{<description>}\end{tablebox}
% Replaces the buggy \caption* (which emitted a stray "טבלה N : *" line) and
% gives every table one identical frame + numbered caption "טבלה N.M.".
% Optional intro text may be placed at the TOP, before the tabular, inside the box.
% ============================================================================
\definecolor{cTblLine}{HTML}{B9C4D2}  \definecolor{cTblBg}{HTML}{FAFBFC}
% v3.23: class-aware — under `article` there is no chapter counter (the promised
% report->article switch used to die with "No counter 'chapter' defined").
\makeatletter
\@ifundefined{c@chapter}{%
  \newcounter{tbl}[section]\renewcommand{\thetbl}{\arabic{tbl}}%
}{%
  \newcounter{tbl}[chapter]\renewcommand{\thetbl}{\thechapter.\arabic{tbl}}%
}
\makeatother
\newmdenv[linecolor=cTblLine,backgroundcolor=cTblBg,linewidth=0.9pt,
  innertopmargin=9pt,innerbottommargin=8pt,skipabove=11pt,skipbelow=11pt,
  innerleftmargin=11pt,innerrightmargin=11pt,roundcorner=4pt]{tablebox}
\newcommand{\tabcap}[1]{\refstepcounter{tbl}\par\addvspace{6pt}%
  {\centering\small\textcolor{cGray}{\textbf{טבלה~\thetbl.}}\enspace #1\par}}

% ============================================================================
% LTR islands.  EVERY tikzpicture / quantikz MUST live inside one, or the whole
% figure's coordinates and Latin labels flip. (Hebrew node text still uses \h{}.)
% ============================================================================
\newenvironment{qcirc}{\par\medskip\begin{otherlanguage}{english}\begin{center}}
  {\end{center}\end{otherlanguage}\par\medskip}
\newenvironment{tikzpic}{\par\medskip\begin{otherlanguage}{english}\begin{center}}
  {\end{center}\end{otherlanguage}\par\medskip}

% ============================================================================
% Code listings.  ★★ v3.4 THE BOX-OVERFLOW FIX:
%   resetmargins=true   → margins measured from the ENCLOSING box, not the page
%                         (without it, a listing inside an mdframed box spills
%                          its frame + line-numbers out past the box border).
%   framexleftmargin=0  → the frame (left bar) does not extend left of the code.
%   xleftmargin keeps a small indent so the line numbers have room INSIDE.
% asmblock wraps the listing in an LTR island (code is LTR).
% ============================================================================
\lstdefinestyle{asm}{
  basicstyle=\ttfamily\small,
  backgroundcolor=\color{cBlueBg!55},
  frame=leftline,framerule=2.5pt,rulecolor=\color{cBlue},
  numbers=left,numberstyle=\tiny\color{cGray},numbersep=8pt,
  xleftmargin=2.2em,framexleftmargin=0pt,framexrightmargin=0pt,resetmargins=true,
  breaklines=true,columns=fullflexible,keepspaces=true,
  showstringspaces=false,aboveskip=8pt,belowskip=8pt,
  commentstyle=\color{cGreen},
}
\lstdefinestyle{cstyle}{
  language=C,basicstyle=\ttfamily\small,
  backgroundcolor=\color{cGreenBg!55},
  frame=leftline,framerule=2.5pt,rulecolor=\color{cGreen},
  numbers=left,numberstyle=\tiny\color{cGray},numbersep=8pt,
  xleftmargin=2.2em,framexleftmargin=0pt,framexrightmargin=0pt,resetmargins=true,
  breaklines=true,columns=fullflexible,keepspaces=true,
  showstringspaces=false,keywordstyle=\color{cBlue}\bfseries,
  commentstyle=\color{cGray}\itshape,aboveskip=8pt,belowskip=8pt,
}
\newenvironment{asmblock}{\par\begin{otherlanguage}{english}}{\end{otherlanguage}\par}

\usepackage[hidelinks,unicode]{hyperref}   % LAST, after colors. unicode = Hebrew bookmarks.

\setcounter{tocdepth}{1}
\setcounter{secnumdepth}{2}
     % (or paste these lines at the top)
%   \begin{document}  ...  \end{document}
%
% Compile with:  lualatex -interaction=nonstopmode doc.tex   (twice, for TOC)
% Requires the Hebrew fonts installed (run scripts/setup_fonts.sh first).
%
% SHORT doc (exam + solution): change `report` -> `article`.
% STUDY doc meant to be read/annotated: 12pt + \linespread{1.3} below are
%   deliberate readability choices (NOT padding) — keep or trim to taste.
% ============================================================================
\documentclass[12pt]{report}          % 12pt: readable for a study guide; 11pt is fine too
\usepackage[no-math]{fontspec}        % (1) ★ no-math — THE critical line (see recipe.md)
\usepackage[bidi=basic]{babel}        % (2) RTL handling built into the engine
\babelprovide[import,main]{hebrew}    % (3) Hebrew = main language, default RTL
\babelprovide[import]{english}        % (4) English secondary, for correct LTR runs
% Latin fallback registered BEFORE \babelfont so a lone English word in Hebrew
% prose does not tofu. Fonts resolved by FAMILY NAME (portable; no hard path).
\directlua{luaotfload.add_fallback("hebfb",{"DejaVuSerif:mode=node;"})}
% v3.22: Frank Ruhl Libre ships NO italic face. Instead of silently-upright
% \emph (v<=3.20) or bold-as-emphasis (v3.21), synthesize a slanted face with
% FakeSlant — real, VISIBLE italic emphasis with correct \emph nesting
% semantics (inner \emph toggles back upright). Verified empirically at
% 300dpi: slant clearly visible, 0 errors, 0 missing glyphs.
\babelfont{rm}[RawFeature={fallback=hebfb},
  ItalicFont={Frank Ruhl Libre}, ItalicFeatures={FakeSlant=0.22},
  BoldItalicFont={Frank Ruhl Libre Bold}, BoldItalicFeatures={FakeSlant=0.22}
  ]{Frank Ruhl Libre}
\babelfont{sf}[RawFeature={fallback=hebfb}]{Heebo}
\babelfont{tt}[RawFeature={fallback=hebfb}]{DejaVu Sans Mono}

\usepackage{amsmath,amssymb,mathtools}
\usepackage[margin=2.3cm,headheight=14pt]{geometry}
\linespread{1.3}                      % readable line spacing for study/annotation docs
\usepackage{xcolor}
\usepackage{mdframed}
\usepackage{listings}                 % code listings (assembly / C / pseudo-code)
\usepackage{tikz}
\usetikzlibrary{arrows.meta,positioning,calc,shapes.geometric,fit,backgrounds}
\usepackage{pgfplots}\pgfplotsset{compat=1.18}
\usepackage{quantikz}                 % quantum circuits (remove if unused)
\usepackage{booktabs}
\usepackage{longtable}
\usepackage{array}
\usepackage{float}                    % [H] = place table exactly here (reference doc: no drifting)
\usepackage{enumitem}
\usepackage{needspace}

% ============================================================================
% ★★ v3.4 — THE DIRECTION-ISOLATION TRIO (the heart of mixing Hebrew + Latin)
% Mental model:
%   \h{...}    Hebrew (RTL) island      — Hebrew text inside an LTR context
%                                          (TikZ nodes, the english `tikzpic`).
%   \en{...}   English (LTR) isolate    — Latin text / parentheses inside a
%                                          Hebrew (RTL) context. KEEPS PARENS
%                                          AND ORDER CORRECT.
%   \code{...} \en + monospace          — code inside Hebrew / inside a frametitle.
% CRITICAL: never wrap Latin in \h{} (it is for HEBREW). Inside a TikZ node,
% Latin in \h{} renders REVERSED (Datapath->htapataD, CPU->UPC) and parens
% mirror. Isolate every Latin/paren chunk with \en{} or \code{}. See
% references/figures-boxes-listings.md.
% ---- TEXT INSIDE MATH (v3.8) ----------------------------------------------
% Put ANY text that sits inside math in a \text{}. The ONLY thing that is
% invalid in math is a bare \h/\en (= \foreignlanguage) used DIRECTLY, outside
% \text -> "\rmfamily invalid in math mode". Inside \text:
%   Hebrew label           -> \text{...}            (bare Hebrew renders fine;
%                              \h is OPTIONAL/harmless: \text{\h{...}} also works
%                              and both render identically + scale in scripts)
%   plain Latin word        -> \text{...}            (no marker needed)
%   Latin WITH ( ) or [ ]   -> \text{\en{...}}       (\en forces LTR; a plain
%                              \text mirrors the parens in RTL -> "Spec (TNR(" )
%   f(Hebrew/text arg)      -> \text{f}(\text{...})  -- keep the parens at MATH
%                              level, NOT \text{f(arg)} (parens inside the box mirror)
%   code in math            -> \mathtt{...}  (or \texttt{}, math-guarded below)
% NEVER write \en{...} or \h{...} straight into $...$ / \[..\] (outside \text);
% always via \text{}. (\mbox{...} also works but does not scale in scripts.)
% ============================================================================
\newcommand{\h}[1]{\foreignlanguage{hebrew}{#1}}
\newcommand{\en}[1]{\foreignlanguage{english}{#1}}

% ============================================================================
% ★★ — fixes for defects that are INVISIBLE in the log and in the checks
% and only surface on a VISUAL RENDER of the PDF. (Each verified empirically.)
% ----------------------------------------------------------------------------
% v3.22: \emph is now REAL slanted emphasis (FakeSlant italic face defined on
% the \babelfont{rm} line above), replacing the v3.21 \emph->\textbf hack: it
% preserves \emph's toggle semantics and keeps bold free for its own role.
% Prefer bold-as-emphasis for your doc? one line: \renewcommand{\emph}[1]{\textbf{#1}}
% page number: force the digits LTR so they never reverse (18 -> 81) in the RTL
%   footer. That reversal is a NON-DETERMINISTIC bidi glitch (observed once in
%   ~80 pp); this LuaLaTeX build is not byte-reproducible, so the defensive
%   LTR wrap is the only reliable guard. Latin digits are always LTR anyway.
\renewcommand{\thepage}{\foreignlanguage{english}{\arabic{page}}}
% emergencystretch: Hebrew does not hyphenate, so an unbreakable LTR island
%   (\en{...} / an inline math run) landing at the margin can leave a small
%   (~2-13pt) overfull with no legal breakpoint. This rescues ONLY those
%   inline-prose tie-breaks; it is a ceiling used only on offending lines and
%   does NOT touch display-math / listing / table overflow (those still surface
%   in the log for the content-level overflow pass). Verified: baseline overfull
%   -> 0 with this set, with no change to pagination or underfull.
\setlength{\emergencystretch}{3em}
% em-dash: --- does NOT become — under Frank Ruhl (the font lacks tlig; verified
%   the width is identical with/without Ligatures=TeX, while the mechanism works
%   on Latin Modern). TYPE THE CHARACTER — (U+2014) DIRECTLY; likewise – (U+2013).
%   check_bidi_figures.py flags a --/--- that is adjacent to Hebrew.
% ============================================================================

% ---- STRUCTURAL DIRECTION-SAFETY for \texttt (v3.7) -------------------------
% A parenthesis, bracket, or $ INSIDE \texttt{} mirrors/migrates in RTL prose
% (MIPS 0($s0)->"(s0$)0", $rs->"rs$", arr[i]->"[arr[i") because \texttt sets the
% font but NOT the direction. Redefining \texttt to LTR-isolate its content fixes
% the WHOLE class at once (parens, brackets, $) so plain \texttt{lw $t0,0($s0)}
% is safe — no need to remember \code{}. Guarded by \ifmmode because
% \foreignlanguage is invalid in math (there, plain monospace is already LTR).
% NB: \texttt{Hebrew} stays tofu — the monospace font has no Hebrew anyway; put
% Hebrew in \h{}/normal text, never in \texttt{}. To disable, delete this block.
\let\origtexttt\texttt
\renewcommand{\texttt}[1]{\ifmmode\origtexttt{#1}\else\foreignlanguage{english}{\origtexttt{#1}}\fi}
\newcommand{\code}[1]{\ifmmode\origtexttt{#1}\else\foreignlanguage{english}{\origtexttt{#1}}\fi}

% ---- a thin decorative rule helper ----
\newcommand{\sep}{\par\vspace{2pt}\noindent\textcolor{cRule}{\rule{\linewidth}{0.6pt}}\par\vspace{4pt}}

% ---- colors ----
\definecolor{cBlue}{HTML}{1F4E79}   \definecolor{cBlueBg}{HTML}{EAF1F8}
\definecolor{cGreen}{HTML}{2E6B4F}  \definecolor{cGreenBg}{HTML}{EAF4EE}
\definecolor{cAmber}{HTML}{9C6A00}  \definecolor{cAmberBg}{HTML}{FBF3E2}
\definecolor{cPurple}{HTML}{5B2C6F} \definecolor{cPurpleBg}{HTML}{F3ECF7}
\definecolor{cTeal}{HTML}{0E6E6E}   \definecolor{cTealBg}{HTML}{E6F4F4}
\definecolor{cRed}{HTML}{8E2B2B}    \definecolor{cRedBg}{HTML}{FBEAEA}
\definecolor{cGray}{HTML}{555555}   \definecolor{cRule}{HTML}{1F4E79}

% ---- math operators ----
\DeclareMathOperator{\Tr}{Tr}
\DeclareMathOperator{\rank}{rank}
\DeclareMathOperator{\diag}{diag}
\DeclareMathOperator{\sgn}{sgn}

% ---- bra-ket (quantikz defines \ket/\bra/\braket; add the rest safely) ----
\providecommand{\ket}[1]{\left|#1\right\rangle}
\providecommand{\bra}[1]{\left\langle#1\right|}
\providecommand{\braket}[2]{\left\langle#1\middle|#2\right\rangle}
\newcommand{\dyad}[1]{\left|#1\right\rangle\!\left\langle#1\right|}
% v3.21: general outer product |a><b| and expectation value <A> — agents reach
% for these by reflex (esp. \expval); quantikz defines neither, so their absence
% threw "Undefined control sequence" in a real build. \providecommand = safe.
\providecommand{\ketbra}[2]{\left|#1\right\rangle\!\left\langle#2\right|}
\providecommand{\expval}[1]{\left\langle#1\right\rangle}
\newcommand{\abs}[1]{\left|#1\right|}
\newcommand{\norm}[1]{\left\lVert#1\right\rVert}
\newcommand{\half}{\tfrac12}
\newcommand{\R}{\mathbb{R}}\newcommand{\C}{\mathbb{C}}\newcommand{\Z}{\mathbb{Z}}
% identity operator. NB: \mathbb{1} silently renders a WRONG glyph (msbm has no
% blackboard digits). Use \mathbb{I} = 𝕀 (works with amssymb, no extra package).
% Do NOT "fix" this back to \mathbb{1}.
\newcommand{\id}{\mathbb{I}}

% ============================================================================
% Colored boxes — mdframed (NOT tcolorbox, which breaks under bidi).
% ★★ v3.4: innerleftmargin/innerrightmargin give code listings room so their
% frame/line-numbers do not spill out the box border (see listing styles below).
% ============================================================================
\newmdenv[linecolor=cBlue,backgroundcolor=cBlueBg,linewidth=1.1pt,
  frametitlebackgroundcolor=cBlue,frametitlefontcolor=white,
  innertopmargin=6pt,innerbottommargin=7pt,skipabove=9pt,skipbelow=9pt,
  innerleftmargin=9pt,innerrightmargin=9pt,roundcorner=3pt]{defbox}   % הגדרה
\newmdenv[linecolor=cGreen,backgroundcolor=cGreenBg,linewidth=1.1pt,
  frametitlebackgroundcolor=cGreen,frametitlefontcolor=white,
  innertopmargin=6pt,innerbottommargin=7pt,skipabove=9pt,skipbelow=9pt,
  innerleftmargin=9pt,innerrightmargin=9pt,roundcorner=3pt]{thmbox}   % משפט/טענה
\newmdenv[linecolor=cAmber,backgroundcolor=cAmberBg,linewidth=1.1pt,
  frametitlebackgroundcolor=cAmber,frametitlefontcolor=white,
  innertopmargin=6pt,innerbottommargin=7pt,skipabove=9pt,skipbelow=9pt,
  innerleftmargin=9pt,innerrightmargin=9pt,roundcorner=3pt]{notebox}  % הערה
\newmdenv[linecolor=cPurple,backgroundcolor=cPurpleBg,linewidth=1.1pt,
  frametitlebackgroundcolor=cPurple,frametitlefontcolor=white,
  innertopmargin=6pt,innerbottommargin=7pt,skipabove=9pt,skipbelow=9pt,
  innerleftmargin=9pt,innerrightmargin=9pt,roundcorner=3pt]{exbox}    % דוגמה
\newmdenv[linecolor=cRed,backgroundcolor=cRedBg,linewidth=1.1pt,
  frametitlebackgroundcolor=cRed,frametitlefontcolor=white,
  innertopmargin=6pt,innerbottommargin=7pt,skipabove=9pt,skipbelow=9pt,
  innerleftmargin=9pt,innerrightmargin=9pt,roundcorner=3pt]{warnbox}  % טעות נפוצה
\newmdenv[linecolor=cTeal,backgroundcolor=cTealBg,linewidth=1.1pt,
  frametitlebackgroundcolor=cTeal,frametitlefontcolor=white,
  innertopmargin=6pt,innerbottommargin=7pt,skipabove=9pt,skipbelow=9pt,
  innerleftmargin=9pt,innerrightmargin=9pt,roundcorner=3pt]{keybox}   % רעיון מרכזי

% ============================================================================
% Unified TABLE presentation: subtle uniform slate frame + auto-numbered
% caption rendered INSIDE the frame. Use:
%   \begin{tablebox}\centering\small  <tabular>  \tabcap{<description>}\end{tablebox}
% Replaces the buggy \caption* (which emitted a stray "טבלה N : *" line) and
% gives every table one identical frame + numbered caption "טבלה N.M.".
% Optional intro text may be placed at the TOP, before the tabular, inside the box.
% ============================================================================
\definecolor{cTblLine}{HTML}{B9C4D2}  \definecolor{cTblBg}{HTML}{FAFBFC}
% v3.23: class-aware — under `article` there is no chapter counter (the promised
% report->article switch used to die with "No counter 'chapter' defined").
\makeatletter
\@ifundefined{c@chapter}{%
  \newcounter{tbl}[section]\renewcommand{\thetbl}{\arabic{tbl}}%
}{%
  \newcounter{tbl}[chapter]\renewcommand{\thetbl}{\thechapter.\arabic{tbl}}%
}
\makeatother
\newmdenv[linecolor=cTblLine,backgroundcolor=cTblBg,linewidth=0.9pt,
  innertopmargin=9pt,innerbottommargin=8pt,skipabove=11pt,skipbelow=11pt,
  innerleftmargin=11pt,innerrightmargin=11pt,roundcorner=4pt]{tablebox}
\newcommand{\tabcap}[1]{\refstepcounter{tbl}\par\addvspace{6pt}%
  {\centering\small\textcolor{cGray}{\textbf{טבלה~\thetbl.}}\enspace #1\par}}

% ============================================================================
% LTR islands.  EVERY tikzpicture / quantikz MUST live inside one, or the whole
% figure's coordinates and Latin labels flip. (Hebrew node text still uses \h{}.)
% ============================================================================
\newenvironment{qcirc}{\par\medskip\begin{otherlanguage}{english}\begin{center}}
  {\end{center}\end{otherlanguage}\par\medskip}
\newenvironment{tikzpic}{\par\medskip\begin{otherlanguage}{english}\begin{center}}
  {\end{center}\end{otherlanguage}\par\medskip}

% ============================================================================
% Code listings.  ★★ v3.4 THE BOX-OVERFLOW FIX:
%   resetmargins=true   → margins measured from the ENCLOSING box, not the page
%                         (without it, a listing inside an mdframed box spills
%                          its frame + line-numbers out past the box border).
%   framexleftmargin=0  → the frame (left bar) does not extend left of the code.
%   xleftmargin keeps a small indent so the line numbers have room INSIDE.
% asmblock wraps the listing in an LTR island (code is LTR).
% ============================================================================
\lstdefinestyle{asm}{
  basicstyle=\ttfamily\small,
  backgroundcolor=\color{cBlueBg!55},
  frame=leftline,framerule=2.5pt,rulecolor=\color{cBlue},
  numbers=left,numberstyle=\tiny\color{cGray},numbersep=8pt,
  xleftmargin=2.2em,framexleftmargin=0pt,framexrightmargin=0pt,resetmargins=true,
  breaklines=true,columns=fullflexible,keepspaces=true,
  showstringspaces=false,aboveskip=8pt,belowskip=8pt,
  commentstyle=\color{cGreen},
}
\lstdefinestyle{cstyle}{
  language=C,basicstyle=\ttfamily\small,
  backgroundcolor=\color{cGreenBg!55},
  frame=leftline,framerule=2.5pt,rulecolor=\color{cGreen},
  numbers=left,numberstyle=\tiny\color{cGray},numbersep=8pt,
  xleftmargin=2.2em,framexleftmargin=0pt,framexrightmargin=0pt,resetmargins=true,
  breaklines=true,columns=fullflexible,keepspaces=true,
  showstringspaces=false,keywordstyle=\color{cBlue}\bfseries,
  commentstyle=\color{cGray}\itshape,aboveskip=8pt,belowskip=8pt,
}
\newenvironment{asmblock}{\par\begin{otherlanguage}{english}}{\end{otherlanguage}\par}

\usepackage[hidelinks,unicode]{hyperref}   % LAST, after colors. unicode = Hebrew bookmarks.

\setcounter{tocdepth}{1}
\setcounter{secnumdepth}{2}
     % (or paste these lines at the top)
%   \begin{document}  ...  \end{document}
%
% Compile with:  lualatex -interaction=nonstopmode doc.tex   (twice, for TOC)
% Requires the Hebrew fonts installed (run scripts/setup_fonts.sh first).
%
% SHORT doc (exam + solution): change `report` -> `article`.
% STUDY doc meant to be read/annotated: 12pt + \linespread{1.3} below are
%   deliberate readability choices (NOT padding) — keep or trim to taste.
% ============================================================================
\documentclass[12pt]{report}          % 12pt: readable for a study guide; 11pt is fine too
\usepackage[no-math]{fontspec}        % (1) ★ no-math — THE critical line (see recipe.md)
\usepackage[bidi=basic]{babel}        % (2) RTL handling built into the engine
\babelprovide[import,main]{hebrew}    % (3) Hebrew = main language, default RTL
\babelprovide[import]{english}        % (4) English secondary, for correct LTR runs
% Latin fallback registered BEFORE \babelfont so a lone English word in Hebrew
% prose does not tofu. Fonts resolved by FAMILY NAME (portable; no hard path).
\directlua{luaotfload.add_fallback("hebfb",{"DejaVuSerif:mode=node;"})}
% v3.22: Frank Ruhl Libre ships NO italic face. Instead of silently-upright
% \emph (v<=3.20) or bold-as-emphasis (v3.21), synthesize a slanted face with
% FakeSlant — real, VISIBLE italic emphasis with correct \emph nesting
% semantics (inner \emph toggles back upright). Verified empirically at
% 300dpi: slant clearly visible, 0 errors, 0 missing glyphs.
\babelfont{rm}[RawFeature={fallback=hebfb},
  ItalicFont={Frank Ruhl Libre}, ItalicFeatures={FakeSlant=0.22},
  BoldItalicFont={Frank Ruhl Libre Bold}, BoldItalicFeatures={FakeSlant=0.22}
  ]{Frank Ruhl Libre}
\babelfont{sf}[RawFeature={fallback=hebfb}]{Heebo}
\babelfont{tt}[RawFeature={fallback=hebfb}]{DejaVu Sans Mono}

\usepackage{amsmath,amssymb,mathtools}
\usepackage[margin=2.3cm,headheight=14pt]{geometry}
\linespread{1.3}                      % readable line spacing for study/annotation docs
\usepackage{xcolor}
\usepackage{mdframed}
\usepackage{listings}                 % code listings (assembly / C / pseudo-code)
\usepackage{tikz}
\usetikzlibrary{arrows.meta,positioning,calc,shapes.geometric,fit,backgrounds}
\usepackage{pgfplots}\pgfplotsset{compat=1.18}
\usepackage{quantikz}                 % quantum circuits (remove if unused)
\usepackage{booktabs}
\usepackage{longtable}
\usepackage{array}
\usepackage{float}                    % [H] = place table exactly here (reference doc: no drifting)
\usepackage{enumitem}
\usepackage{needspace}

% ============================================================================
% ★★ v3.4 — THE DIRECTION-ISOLATION TRIO (the heart of mixing Hebrew + Latin)
% Mental model:
%   \h{...}    Hebrew (RTL) island      — Hebrew text inside an LTR context
%                                          (TikZ nodes, the english `tikzpic`).
%   \en{...}   English (LTR) isolate    — Latin text / parentheses inside a
%                                          Hebrew (RTL) context. KEEPS PARENS
%                                          AND ORDER CORRECT.
%   \code{...} \en + monospace          — code inside Hebrew / inside a frametitle.
% CRITICAL: never wrap Latin in \h{} (it is for HEBREW). Inside a TikZ node,
% Latin in \h{} renders REVERSED (Datapath->htapataD, CPU->UPC) and parens
% mirror. Isolate every Latin/paren chunk with \en{} or \code{}. See
% references/figures-boxes-listings.md.
% ---- TEXT INSIDE MATH (v3.8) ----------------------------------------------
% Put ANY text that sits inside math in a \text{}. The ONLY thing that is
% invalid in math is a bare \h/\en (= \foreignlanguage) used DIRECTLY, outside
% \text -> "\rmfamily invalid in math mode". Inside \text:
%   Hebrew label           -> \text{...}            (bare Hebrew renders fine;
%                              \h is OPTIONAL/harmless: \text{\h{...}} also works
%                              and both render identically + scale in scripts)
%   plain Latin word        -> \text{...}            (no marker needed)
%   Latin WITH ( ) or [ ]   -> \text{\en{...}}       (\en forces LTR; a plain
%                              \text mirrors the parens in RTL -> "Spec (TNR(" )
%   f(Hebrew/text arg)      -> \text{f}(\text{...})  -- keep the parens at MATH
%                              level, NOT \text{f(arg)} (parens inside the box mirror)
%   code in math            -> \mathtt{...}  (or \texttt{}, math-guarded below)
% NEVER write \en{...} or \h{...} straight into $...$ / \[..\] (outside \text);
% always via \text{}. (\mbox{...} also works but does not scale in scripts.)
% ============================================================================
\newcommand{\h}[1]{\foreignlanguage{hebrew}{#1}}
\newcommand{\en}[1]{\foreignlanguage{english}{#1}}

% ============================================================================
% ★★ — fixes for defects that are INVISIBLE in the log and in the checks
% and only surface on a VISUAL RENDER of the PDF. (Each verified empirically.)
% ----------------------------------------------------------------------------
% v3.22: \emph is now REAL slanted emphasis (FakeSlant italic face defined on
% the \babelfont{rm} line above), replacing the v3.21 \emph->\textbf hack: it
% preserves \emph's toggle semantics and keeps bold free for its own role.
% Prefer bold-as-emphasis for your doc? one line: \renewcommand{\emph}[1]{\textbf{#1}}
% page number: force the digits LTR so they never reverse (18 -> 81) in the RTL
%   footer. That reversal is a NON-DETERMINISTIC bidi glitch (observed once in
%   ~80 pp); this LuaLaTeX build is not byte-reproducible, so the defensive
%   LTR wrap is the only reliable guard. Latin digits are always LTR anyway.
\renewcommand{\thepage}{\foreignlanguage{english}{\arabic{page}}}
% emergencystretch: Hebrew does not hyphenate, so an unbreakable LTR island
%   (\en{...} / an inline math run) landing at the margin can leave a small
%   (~2-13pt) overfull with no legal breakpoint. This rescues ONLY those
%   inline-prose tie-breaks; it is a ceiling used only on offending lines and
%   does NOT touch display-math / listing / table overflow (those still surface
%   in the log for the content-level overflow pass). Verified: baseline overfull
%   -> 0 with this set, with no change to pagination or underfull.
\setlength{\emergencystretch}{3em}
% em-dash: --- does NOT become — under Frank Ruhl (the font lacks tlig; verified
%   the width is identical with/without Ligatures=TeX, while the mechanism works
%   on Latin Modern). TYPE THE CHARACTER — (U+2014) DIRECTLY; likewise – (U+2013).
%   check_bidi_figures.py flags a --/--- that is adjacent to Hebrew.
% ============================================================================

% ---- STRUCTURAL DIRECTION-SAFETY for \texttt (v3.7) -------------------------
% A parenthesis, bracket, or $ INSIDE \texttt{} mirrors/migrates in RTL prose
% (MIPS 0($s0)->"(s0$)0", $rs->"rs$", arr[i]->"[arr[i") because \texttt sets the
% font but NOT the direction. Redefining \texttt to LTR-isolate its content fixes
% the WHOLE class at once (parens, brackets, $) so plain \texttt{lw $t0,0($s0)}
% is safe — no need to remember \code{}. Guarded by \ifmmode because
% \foreignlanguage is invalid in math (there, plain monospace is already LTR).
% NB: \texttt{Hebrew} stays tofu — the monospace font has no Hebrew anyway; put
% Hebrew in \h{}/normal text, never in \texttt{}. To disable, delete this block.
\let\origtexttt\texttt
\renewcommand{\texttt}[1]{\ifmmode\origtexttt{#1}\else\foreignlanguage{english}{\origtexttt{#1}}\fi}
\newcommand{\code}[1]{\ifmmode\origtexttt{#1}\else\foreignlanguage{english}{\origtexttt{#1}}\fi}

% ---- a thin decorative rule helper ----
\newcommand{\sep}{\par\vspace{2pt}\noindent\textcolor{cRule}{\rule{\linewidth}{0.6pt}}\par\vspace{4pt}}

% ---- colors ----
\definecolor{cBlue}{HTML}{1F4E79}   \definecolor{cBlueBg}{HTML}{EAF1F8}
\definecolor{cGreen}{HTML}{2E6B4F}  \definecolor{cGreenBg}{HTML}{EAF4EE}
\definecolor{cAmber}{HTML}{9C6A00}  \definecolor{cAmberBg}{HTML}{FBF3E2}
\definecolor{cPurple}{HTML}{5B2C6F} \definecolor{cPurpleBg}{HTML}{F3ECF7}
\definecolor{cTeal}{HTML}{0E6E6E}   \definecolor{cTealBg}{HTML}{E6F4F4}
\definecolor{cRed}{HTML}{8E2B2B}    \definecolor{cRedBg}{HTML}{FBEAEA}
\definecolor{cGray}{HTML}{555555}   \definecolor{cRule}{HTML}{1F4E79}

% ---- math operators ----
\DeclareMathOperator{\Tr}{Tr}
\DeclareMathOperator{\rank}{rank}
\DeclareMathOperator{\diag}{diag}
\DeclareMathOperator{\sgn}{sgn}

% ---- bra-ket (quantikz defines \ket/\bra/\braket; add the rest safely) ----
\providecommand{\ket}[1]{\left|#1\right\rangle}
\providecommand{\bra}[1]{\left\langle#1\right|}
\providecommand{\braket}[2]{\left\langle#1\middle|#2\right\rangle}
\newcommand{\dyad}[1]{\left|#1\right\rangle\!\left\langle#1\right|}
% v3.21: general outer product |a><b| and expectation value <A> — agents reach
% for these by reflex (esp. \expval); quantikz defines neither, so their absence
% threw "Undefined control sequence" in a real build. \providecommand = safe.
\providecommand{\ketbra}[2]{\left|#1\right\rangle\!\left\langle#2\right|}
\providecommand{\expval}[1]{\left\langle#1\right\rangle}
\newcommand{\abs}[1]{\left|#1\right|}
\newcommand{\norm}[1]{\left\lVert#1\right\rVert}
\newcommand{\half}{\tfrac12}
\newcommand{\R}{\mathbb{R}}\newcommand{\C}{\mathbb{C}}\newcommand{\Z}{\mathbb{Z}}
% identity operator. NB: \mathbb{1} silently renders a WRONG glyph (msbm has no
% blackboard digits). Use \mathbb{I} = 𝕀 (works with amssymb, no extra package).
% Do NOT "fix" this back to \mathbb{1}.
\newcommand{\id}{\mathbb{I}}

% ============================================================================
% Colored boxes — mdframed (NOT tcolorbox, which breaks under bidi).
% ★★ v3.4: innerleftmargin/innerrightmargin give code listings room so their
% frame/line-numbers do not spill out the box border (see listing styles below).
% ============================================================================
\newmdenv[linecolor=cBlue,backgroundcolor=cBlueBg,linewidth=1.1pt,
  frametitlebackgroundcolor=cBlue,frametitlefontcolor=white,
  innertopmargin=6pt,innerbottommargin=7pt,skipabove=9pt,skipbelow=9pt,
  innerleftmargin=9pt,innerrightmargin=9pt,roundcorner=3pt]{defbox}   % הגדרה
\newmdenv[linecolor=cGreen,backgroundcolor=cGreenBg,linewidth=1.1pt,
  frametitlebackgroundcolor=cGreen,frametitlefontcolor=white,
  innertopmargin=6pt,innerbottommargin=7pt,skipabove=9pt,skipbelow=9pt,
  innerleftmargin=9pt,innerrightmargin=9pt,roundcorner=3pt]{thmbox}   % משפט/טענה
\newmdenv[linecolor=cAmber,backgroundcolor=cAmberBg,linewidth=1.1pt,
  frametitlebackgroundcolor=cAmber,frametitlefontcolor=white,
  innertopmargin=6pt,innerbottommargin=7pt,skipabove=9pt,skipbelow=9pt,
  innerleftmargin=9pt,innerrightmargin=9pt,roundcorner=3pt]{notebox}  % הערה
\newmdenv[linecolor=cPurple,backgroundcolor=cPurpleBg,linewidth=1.1pt,
  frametitlebackgroundcolor=cPurple,frametitlefontcolor=white,
  innertopmargin=6pt,innerbottommargin=7pt,skipabove=9pt,skipbelow=9pt,
  innerleftmargin=9pt,innerrightmargin=9pt,roundcorner=3pt]{exbox}    % דוגמה
\newmdenv[linecolor=cRed,backgroundcolor=cRedBg,linewidth=1.1pt,
  frametitlebackgroundcolor=cRed,frametitlefontcolor=white,
  innertopmargin=6pt,innerbottommargin=7pt,skipabove=9pt,skipbelow=9pt,
  innerleftmargin=9pt,innerrightmargin=9pt,roundcorner=3pt]{warnbox}  % טעות נפוצה
\newmdenv[linecolor=cTeal,backgroundcolor=cTealBg,linewidth=1.1pt,
  frametitlebackgroundcolor=cTeal,frametitlefontcolor=white,
  innertopmargin=6pt,innerbottommargin=7pt,skipabove=9pt,skipbelow=9pt,
  innerleftmargin=9pt,innerrightmargin=9pt,roundcorner=3pt]{keybox}   % רעיון מרכזי

% ============================================================================
% Unified TABLE presentation: subtle uniform slate frame + auto-numbered
% caption rendered INSIDE the frame. Use:
%   \begin{tablebox}\centering\small  <tabular>  \tabcap{<description>}\end{tablebox}
% Replaces the buggy \caption* (which emitted a stray "טבלה N : *" line) and
% gives every table one identical frame + numbered caption "טבלה N.M.".
% Optional intro text may be placed at the TOP, before the tabular, inside the box.
% ============================================================================
\definecolor{cTblLine}{HTML}{B9C4D2}  \definecolor{cTblBg}{HTML}{FAFBFC}
% v3.23: class-aware — under `article` there is no chapter counter (the promised
% report->article switch used to die with "No counter 'chapter' defined").
\makeatletter
\@ifundefined{c@chapter}{%
  \newcounter{tbl}[section]\renewcommand{\thetbl}{\arabic{tbl}}%
}{%
  \newcounter{tbl}[chapter]\renewcommand{\thetbl}{\thechapter.\arabic{tbl}}%
}
\makeatother
\newmdenv[linecolor=cTblLine,backgroundcolor=cTblBg,linewidth=0.9pt,
  innertopmargin=9pt,innerbottommargin=8pt,skipabove=11pt,skipbelow=11pt,
  innerleftmargin=11pt,innerrightmargin=11pt,roundcorner=4pt]{tablebox}
\newcommand{\tabcap}[1]{\refstepcounter{tbl}\par\addvspace{6pt}%
  {\centering\small\textcolor{cGray}{\textbf{טבלה~\thetbl.}}\enspace #1\par}}

% ============================================================================
% LTR islands.  EVERY tikzpicture / quantikz MUST live inside one, or the whole
% figure's coordinates and Latin labels flip. (Hebrew node text still uses \h{}.)
% ============================================================================
\newenvironment{qcirc}{\par\medskip\begin{otherlanguage}{english}\begin{center}}
  {\end{center}\end{otherlanguage}\par\medskip}
\newenvironment{tikzpic}{\par\medskip\begin{otherlanguage}{english}\begin{center}}
  {\end{center}\end{otherlanguage}\par\medskip}

% ============================================================================
% Code listings.  ★★ v3.4 THE BOX-OVERFLOW FIX:
%   resetmargins=true   → margins measured from the ENCLOSING box, not the page
%                         (without it, a listing inside an mdframed box spills
%                          its frame + line-numbers out past the box border).
%   framexleftmargin=0  → the frame (left bar) does not extend left of the code.
%   xleftmargin keeps a small indent so the line numbers have room INSIDE.
% asmblock wraps the listing in an LTR island (code is LTR).
% ============================================================================
\lstdefinestyle{asm}{
  basicstyle=\ttfamily\small,
  backgroundcolor=\color{cBlueBg!55},
  frame=leftline,framerule=2.5pt,rulecolor=\color{cBlue},
  numbers=left,numberstyle=\tiny\color{cGray},numbersep=8pt,
  xleftmargin=2.2em,framexleftmargin=0pt,framexrightmargin=0pt,resetmargins=true,
  breaklines=true,columns=fullflexible,keepspaces=true,
  showstringspaces=false,aboveskip=8pt,belowskip=8pt,
  commentstyle=\color{cGreen},
}
\lstdefinestyle{cstyle}{
  language=C,basicstyle=\ttfamily\small,
  backgroundcolor=\color{cGreenBg!55},
  frame=leftline,framerule=2.5pt,rulecolor=\color{cGreen},
  numbers=left,numberstyle=\tiny\color{cGray},numbersep=8pt,
  xleftmargin=2.2em,framexleftmargin=0pt,framexrightmargin=0pt,resetmargins=true,
  breaklines=true,columns=fullflexible,keepspaces=true,
  showstringspaces=false,keywordstyle=\color{cBlue}\bfseries,
  commentstyle=\color{cGray}\itshape,aboveskip=8pt,belowskip=8pt,
}
\newenvironment{asmblock}{\par\begin{otherlanguage}{english}}{\end{otherlanguage}\par}

\usepackage[hidelinks,unicode]{hyperref}   % LAST, after colors. unicode = Hebrew bookmarks.

\setcounter{tocdepth}{1}
\setcounter{secnumdepth}{2}

\begin{document}

% ---------------------------------------------------------------------------
% Density ruler — one page that *is* the target balance, labeled.
% ---------------------------------------------------------------------------
\chapter*{סרגל צפיפות — היעד הוויזואלי}
\addcontentsline{toc}{chapter}{סרגל צפיפות}
\setcounter{chapter}{1}

\noindent זהו עמוד־כיול. ספור: \textbf{קופסה אחת}, \textbf{שתי משוואות ממוספרות},
\textbf{טבלה אפורה אחת}, והשאר פרוזה. זה היחס ש־\texttt{content-style.md} דורש —
לא קיר קופסאות.

\begin{defbox}[frametitle={הגדרה — זמן גישה ממוצע}]
$\mathrm{AMAT}=T_{\mathrm{hit}}+r_{\mathrm{miss}}T_{\mathrm{miss}}$.
\end{defbox}

המדד מאחד שני מספרים שקל להתעלם מאחד מהם: זמן פגיעה ושיעור החמצה.

\begin{equation}
\label{eq:amat-ruler}
T_{\mathrm{avg}}=T_h+r_m T_m
\end{equation}

אם רק מקטינים את $T_h$ ומשאירים $r_m$ גבוה —~\eqref{eq:amat-ruler} כמעט לא זז.

\begin{equation}
\label{eq:amat-sens}
\frac{\partial T_{\mathrm{avg}}}{\partial r_m}=T_m
\end{equation}

~\eqref{eq:amat-sens} אומרת: רגישות לשיעור החמצה היא בדיוק מחיר ההחמצה.

\begin{tablebox}\centering\small
\begin{tabular}{@{}lll@{}}
\toprule
\textbf{מצב} & \textbf{סדר גודל} & \textbf{פעולה}\\
\midrule
תקין & $\sim 1\times$ & שמור מדידה\\
חשוד & $1.5$–$2\times$ & בודקים תלות\\
שבור & $\ge 2\times$ & תיקון מבני\\
\bottomrule
\end{tabular}
\tabcap{כיול מהיר ל־\en{AMAT}.}
\end{tablebox}

\noindent\textbf{לקח מהעמוד:} רוב השטח הוא פרוזה + משוואות + טבלה אפורה;
הקופסה הצבעונית היא הדגשה, לא שלד המסמך.

% ---------------------------------------------------------------------------
% Chapter 1 — theorem box + short prose proof (no-cloning DNA)
% ---------------------------------------------------------------------------
\chapter{משפט אי־השכפול}
\label{ch:noclone}

\begin{thmbox}[frametitle={\texorpdfstring{$\mathrm{No\text{-}Cloning}$}{No-Cloning}}]
לא קיים אופרטור אוניטרי $U$ המייצר עותק זהה של מצב קוונטי \emph{שרירותי}, כלומר כך ש־
$U(\ket\psi\otimes\ket0)=\ket\psi\otimes\ket\psi$ לכל $\ket\psi$.
\end{thmbox}
\noindent\textbf{הוכחה (בקצרה).} נניח שקיים $U$ כזה עבור שני מצבים $\ket\psi,\ket\phi$. אז
$\braket{\psi}{\phi}=\bra\psi\bra0\,U^\dagger U\,\ket\phi\ket0=\braket{\psi}{\phi}^2$, ולכן
$\braket{\psi}{\phi}\in\{0,1\}$ — כלומר המצבים זהים או אורתוגונליים. שכפול אוניברסלי של מצב
שרירותי סותר את הליניאריות. $\qquad\blacksquare$

% ---------------------------------------------------------------------------
% Chapter 2 — goal sentence + enumerate + circuit (dense coding DNA)
% ---------------------------------------------------------------------------
\chapter{קידוד צפוף}
\label{ch:dense}

המטרה: להעביר \emph{שני} סיביות קלאסיות בעזרת שליחת \emph{קיוביט אחד}, תוך ניצול שזירה משותפת.

\begin{enumerate}\itemsep3pt
\item[(\textbf{א})] אליס ובוב חולקים מראש את
$\ket{\psi_{00}}=\tfrac1{\sqrt2}(\ket{00}+\ket{11})$.
\item[(\textbf{ב})] לפי שני הסיביות, אליס מפעילה על הקיוביט שלה אחד מ־
$\{\,I\to00,\ X\to01,\ Z\to10,\ XZ\to11\,\}$, והמצב המשותף הופך לאחד מארבעת מצבי בל.
\item[(\textbf{ג})] אליס שולחת את הקיוביט; בוב מפעיל $\mathrm{CNOT}$ ואז $H$ ומודד —
ומקבל בדיוק את שני הסיביות.
\end{enumerate}

\begin{qcirc}
\begin{quantikz}
\lstick[2]{$\ket{\psi_{00}}$} & \gate{U\in\{I,X,Z,XZ\}} & \ctrl{1} & \gate{H} & \meter{} \\
 & \qw & \targ{} & \qw & \meter{}
\end{quantikz}
\end{qcirc}
\begin{center}\small קידוד צפוף: אליס מקודדת, שולחת, ובוב מפענח בבסיס בל.\end{center}

% ---------------------------------------------------------------------------
% Chapter 3 — systems: def + numbered eq + short intuition + steps
% ---------------------------------------------------------------------------
\chapter{סכסוך נתונים בצנרת}
\label{ch:hazard}

היחידה נסגרת סביב \en{RAW}: מתי נוצר, איך מודדים אותו, ומה נשבר כשמתעלמים ממנו.

\begin{warnbox}[frametitle={מלכודת — קידום לא תמיד מספיק}]
גם עם קידום מלא, \code{lw} ואחריו שימוש דורשים \en{stall} של מחזור.
\end{warnbox}

\begin{equation}
\label{eq:stall}
T_{\mathrm{stall}}=N_{\mathrm{load\text{-}use}}
\end{equation}

המשמעות של~\eqref{eq:stall}: מספר זוגות \en{load-use} הוא המדד — לא ״יש/אין סכסוך״.

סדר עבודה:
\begin{enumerate}\itemsep3pt
\item[(\textbf{א})] ספור זוגות \en{load-use} לפי~\eqref{eq:stall}.
\item[(\textbf{ב})] בדוק אם קידום מכסה את השאר.
\item[(\textbf{ג})] רק אז שנה תזמון או סדר הוראות.
\end{enumerate}

\end{document}
